%%============================================================================
%% CARTA DE FUNDAÇÃO — Grupo de Usuários Linux e Software Livre -  Estatística e Ciência de Dados
%%                     Universidade Federal do Paraná (UFPR)
%%============================================================================
%% Documento institucional — versão formal
%% Compilar com: pdflatex carta-fundacao.tex  (ou xelatex / lualatex)
%%============================================================================

\documentclass[
  12pt,          % tamanho base da fonte
  a4paper,       % formato de papel
  oneside,       % impressão frente única
  brazil         % idioma principal: português do Brasil
]{article}

%---------------------------------------------------------------------------
% Pacotes essenciais
%---------------------------------------------------------------------------
\usepackage[T1]{fontenc}            % codificação de fonte
\usepackage[utf8]{inputenc}         % entrada UTF-8 (pdflatex)
\usepackage{lmodern}                % fontes Latin Modern (vetoriais)
\usepackage{microtype}              % microtipografia (melhora legibilidade)

\usepackage[brazil]{babel}          % hifenização e nomes de seções em PT-BR

\usepackage[
  margin=2.5cm,
  top=3cm,
  bottom=3cm,
  headheight=25pt
]{geometry}                         % margens do documento

\usepackage{setspace}               % controle de espaçamento entre linhas
\usepackage{xcolor}                 % suporte a cores
\usepackage{titlesec}               % formatação de títulos de seções
\usepackage{titling}                % formatação do título do documento
\usepackage{fancyhdr}               % cabeçalhos e rodapés personalizados
\usepackage{enumitem}               % personalização de listas
\usepackage{lettrine}               % capitulares (letra inicial decorada)
\usepackage{graphicx}               % suporte a imagens (se necessário)
\usepackage{hyperref}               % hiperlinks e metadados do PDF
\usepackage{bookmark}               % bookmarks do PDF mais limpos

%---------------------------------------------------------------------------
% Cores institucionais
%---------------------------------------------------------------------------
\definecolor{gulecdBlue}{HTML}{1B3A5C}   % azul escuro institucional
\definecolor{gulecdAccent}{HTML}{2E86AB} % azul de destaque
\definecolor{gulecdGray}{HTML}{555555}   % cinza para texto secundário
\definecolor{gulecdLight}{HTML}{F5F7FA}  % fundo claro para caixas
\definecolor{gulecdRule}{HTML}{D0D5DD}   % cor de filetes e réguas

%---------------------------------------------------------------------------
% Metadados do PDF
%---------------------------------------------------------------------------
\hypersetup{
  pdftitle    = {Carta de Fundação — GULECD/UFPR},
  pdfauthor   = {Grupo de Usuários Linux e Software Livre - Estatística e Ciência de Dados UFPR},
  pdfsubject  = {Carta de Fundação institucional do GULECD},
  pdfkeywords = {carta de fundação, software livre, código aberto, ciência de dados, UFPR},
  colorlinks  = true,
  linkcolor   = gulecdAccent,
  urlcolor    = gulecdAccent,
  citecolor   = gulecdAccent
}

%---------------------------------------------------------------------------
% Configurações tipográficas
%---------------------------------------------------------------------------
\onehalfspacing                     % espaçamento entre linhas = 1,5
\setlength{\parindent}{1.5cm}       % recuo de parágrafo
\setlength{\parskip}{0.5em}         % espaço entre parágrafos

%---------------------------------------------------------------------------
% Formato do título do documento
%---------------------------------------------------------------------------
\pretitle{%
  \begin{center}
  \vspace*{1cm}
  \rule{\textwidth}{2pt}\vspace{0.3cm}\\
  \LARGE\bfseries\color{gulecdBlue}
}
\posttitle{%
  \\\vspace{0.3cm}
  \rule{\textwidth}{2pt}
  \end{center}
  \vspace{0.5cm}
}
\preauthor{\begin{center}\large\color{gulecdGray}}
\postauthor{\end{center}}
\predate{\begin{center}\small\color{gulecdGray}}
\postdate{\end{center}}

%---------------------------------------------------------------------------
% Formato das seções
%---------------------------------------------------------------------------
\titleformat{\section}
  [block]
  {\normalfont\Large\bfseries\color{gulecdBlue}}
  {\thesection.}{0.8em}{}
  [\vspace{0.2em}{\color{gulecdRule}\rule{\linewidth}{0.6pt}}]

\titleformat{\subsection}
  [block]
  {\normalfont\large\bfseries\color{gulecdAccent}}
  {\thesubsection.}{0.8em}{}

\titlespacing*{\section}
  {0pt}{2.5ex plus 1ex minus .2ex}{1.5ex plus .2ex}
\titlespacing*{\subsection}
  {0pt}{2ex plus 1ex minus .2ex}{1ex plus .2ex}

%---------------------------------------------------------------------------
% Cabeçalho e rodapé
%---------------------------------------------------------------------------
\pagestyle{fancy}
\fancyhf{}
\fancyhead[L]{\small\color{gulecdGray}\textsc{GULECD/UFPR}}
\fancyhead[R]{\small\color{gulecdGray}\textit{Carta de Fundação}}
\fancyfoot[C]{\small\color{gulecdGray}---\ \thepage\ ---}
\renewcommand{\headrulewidth}{0.4pt}
\renewcommand{\headrule}{%
  \hbox to\headwidth{%
    \color{gulecdRule}\leaders\hrule height \headrulewidth\hfill}}

% Remove cabeçalho/rodapé da primeira página (será redefinido manualmente)
\fancypagestyle{plain}{%
  \fancyhf{}%
  \fancyfoot[C]{\small\color{gulecdGray}---\ \thepage\ ---}%
  \renewcommand{\headrulewidth}{0pt}%
}

%---------------------------------------------------------------------------
% Ambiente personalizado para listas
%---------------------------------------------------------------------------
\newlist{conduta}{itemize}{1}
\setlist[conduta]{
  label      = {\color{gulecdAccent}\textbullet},
  leftmargin = 1.5em,
  itemsep    = 0.3em,
  topsep     = 0.4em
}

%---------------------------------------------------------------------------
% Informações do documento
%---------------------------------------------------------------------------
\title{Carta de Fundação\\[0.3em]
  \Large\mdseries\color{gulecdGray}%
  Grupo de Usuários Linux e Software Livre\\Estatística e Ciência de Dados\\Universidade Federal do Paraná (UFPR)}
%\author{Grupo de Usuários Linux e Software Livre\\Curso de Estatística e Ciência de Dados --- UFPR}
\date{\today}



%============================================================================
% INÍCIO DO DOCUMENTO
%============================================================================
\begin{document}

%--- Capa / Título principal -------------------------------------------------
\thispagestyle{empty}
\maketitle

\newpage

%--- Seção 1: Introdução ----------------------------------------------------
\section{Introdução}

\lettrine[lines=3, lhang=0.1, nindent=0em, slope=-0.5em]{O}{ presente} documento formaliza a criação do Grupo de Usuários Linux e Software Livre do curso de  
Estatística e Ciência de Dados da Universidade Federal do Paraná (UFPR), constituído como uma iniciativa acadêmica e técnica voltada à promoção do 
uso de sistemas baseados em Linux, \textit{software} livre e práticas abertas no contexto da estatística, ciência de dados e computação científica.

Este grupo emerge da convergência entre a necessidade contemporânea de rigor metodológico, reprodutibilidade científica e autonomia tecnológica, e a 
maturidade do ecossistema de \textit{software} livre como infraestrutura essencial para pesquisa, ensino e inovação. Nesse sentido, estabelece-se como um 
espaço de colaboração, formação técnica e produção científica, orientado por princípios formais expressos em seu Manifesto e regulado por seu Código de Conduta.

%--- Seção 2: Justificativa -------------------------------------------------
\section{Justificativa}

A Estatística e Ciência de Dados moderna depende intrinsecamente de ambientes computacionais complexos, frequentemente sustentados por \textit{software} 
de código aberto. Sistemas operacionais baseados em Linux constituem a espinha dorsal de infraestruturas científicas, incluindo computação de alto desempenho, 
\textit{pipelines} de análise de dados, aprendizado de máquina e engenharia de dados.

Entretanto, observa-se uma lacuna entre o uso instrumental dessas tecnologias e sua compreensão aprofundada, 
bem como uma dependência significativa de ferramentas proprietárias que limitam a auditabilidade, a reprodutibilidade e a autonomia de estudantes, professores
pesquisadores e profissionais.

Diante desse cenário, a criação deste grupo justifica-se como um mecanismo de capacitação técnica, organização comunitária e promoção de 
práticas científicas alinhadas aos princípios do conhecimento aberto, visando formar profissionais capazes de operar, compreender e desenvolver 
soluções computacionais robustas, transparentes e eticamente fundamentadas.

%--- Seção 3: Objetivos -----------------------------------------------------
\section{Objetivos}

\subsection{Objetivo Geral}

Promover o uso, estudo e desenvolvimento de tecnologias baseadas em Linux e \textit{software} livre no âmbito da Estatística e Ciência de Dados, 
fomentando a excelência técnica, a colaboração aberta e a produção científica reprodutível.

\subsection{Objetivos Específicos}

\begin{conduta}
  \item Incentivar a adoção de sistemas Linux e outros sistemas operacionais de código aberto como ambiente primário de desenvolvimento científico
  \item Promover o uso de ferramentas de código aberto em análise de dados, engenharia de dados, modelagem estatística e aprendizado de máquina
  \item Capacitar membros em fundamentos computacionais, incluindo sistemas operacionais, redes, programação e infraestrutura
  \item Estimular a produção e compartilhamento de código, dados e metodologias de forma aberta
  \item Organizar eventos técnicos, \textit{workshops}, grupos de estudo e projetos colaborativos
  \item Fomentar práticas de reprodutibilidade científica e versionamento de experimentos
  \item Estabelecer conexões com a comunidade \textit{open source} e científica nacional e internacional
  \item Contribuir para uma infraestrutura digital própria para o Curso de Estatística e Ciência de Dados da UFPR baseada em Linux e Software Livre
\end{conduta}

%--- Seção 4: Princípios Fundamentais ---------------------------------------
\section{Princípios Fundamentais}

O grupo orienta-se pelos princípios estabelecidos em seu Manifesto, que define os valores centrais relacionados à ciência aberta, \textit{software} livre, 
colaboração e rigor científico. Todos os membros comprometem-se a atuar em conformidade com tais princípios.

Adicionalmente, o funcionamento do grupo é regido por seu Código de Conduta, que estabelece normas de convivência, padrões de comportamento e 
mecanismos de mediação de conflitos, assegurando um ambiente respeitoso, inclusivo e produtivo.

%--- Seção 5: Regras Estatutárias -------------------------------------------
\section{Regras Estatutárias}

\subsection{Natureza e Estrutura}

O grupo constitui-se como uma organização acadêmica, sem fins lucrativos, de caráter colaborativo e voluntário, vinculada à comunidade 
discente do curso de Estatística e Ciência de Dados da Universidade Federal do Paraná (UFPR). 
No exercício de suas atividades, o grupo buscará atuar de forma integrada ao ecossistema institucional do curso, promovendo diálogo, 
cooperação e sinergia com outras estruturas discentes e docentes, tais como: o Centro Acadêmico de Estatística (CAEST), o grupo PET 
(Programa de Educação Tutorial) do curso de Estatística e Ciência de Dados da UFPR, 
a coordenação do Curso de Estatística e Ciência de Dados da UFPR e o Departamento de Estatística da UFPR. Estas interações visam o 
fortalecimento das atividades acadêmicas, a organização de eventos, a formação técnica dos estudantes 
e a ampliação do impacto científico e tecnológico no âmbito do curso. Tais cooperações deverão ocorrer de maneira autônoma, 
respeitosa e institucionalmente alinhada, preservando a independência do grupo ao mesmo tempo em que potencializa sua inserção 
e relevância na comunidade universitária.

\subsection{Membros}

São considerados membros todos os indivíduos que participam ativamente dos canais e atividades do grupo e que aceitam, 
explícita ou implicitamente, os termos desta Carta de Fundação, do Manifesto e do Código de Conduta.

Não haverá distinção hierárquica rígida entre os membros no que se refere à produção técnica, 
à participação em discussões e à contribuição científica, sendo incentivada uma estrutura horizontal baseada em mérito, 
contribuição efetiva e responsabilidade compartilhada. Todavia, para fins de organização, continuidade institucional e 
execução coordenada das atividades, o grupo contará com uma Coordenação, composta por membros eleitos através de processo eleitoral periódico.

A Coordenação será escolhida anualmente, mediante eleição conduzida de acordo com regras, critérios e procedimentos previamente definidos no 
Regulamento Eleitoral do grupo, assegurando transparência, participação ampla e igualdade de condições entre os candidatos. 
Poderão candidatar-se membros ativos do grupo, sendo esperado histórico de contribuição, comprometimento com os princípios estabelecidos 
no Manifesto e adesão ao Código de Conduta.

Compete à Coordenação a condução geral do grupo, incluindo planejamento e execução de atividades, organização de eventos, 
gerenciamento dos canais e infraestrutura, mediação de questões administrativas, representação institucional e apoio à coordenação de projetos. 
A atuação da Coordenação deverá observar os princípios de transparência, prestação de contas e alinhamento com os 
valores do grupo, não implicando privilégio executivo, técnico ou científico sobre os demais membros, mas sim responsabilidade funcional e organizacional.


Também compete à Coordenação estimular e apoiar os membros a desenvolver ações de maneira autônoma no contexto do grupo, provendo adequado ambiente
institucional e demais tipos de suporte, mantendo claro que não é privilégio ou direito exclusivo da Coordenação a iniciativa e execução de atividades e projetos. 


Tal participação constitui fundamento pétreo do grupo, assegurando sua continuidade e relevância. 

\subsection{Organização e Governança}

A organização do grupo será realizada por membros voluntários que assumam funções de coordenação, moderação e suporte técnico, incluindo:

\begin{conduta}
  \item Organização de projetos, atividades e eventos
  \item Mediação de discussões e aplicação do Código de Conduta
  \item Manutenção de infraestrutura digital (repositórios, canais, documentação)
  \item Realização de eleições da Coordenação
  \item Representação institucional do grupo
  \item Busca, manutenção e expansão de recursos de qualquer natureza com o objetivo de fomentar, sustentar e desenvolver as ações do grupo 
  em acordo com os princípios do Manifesto, os objetivos descritos neste documento e em consonância com as regras de seu Código de Conduta
\end{conduta}

Decisões relevantes poderão ser tomadas de forma coletiva, priorizando consenso e transparência.

\subsection{Atividades}

As atividades do grupo poderão incluir, mas não se limitam a:

\begin{conduta}
  \item Grupos de estudo e sessões técnicas
  \item Desenvolvimento de projetos colaborativos
  \item \textit{Workshops}, palestras e minicursos
  \item Produção de materiais didáticos e científicos
  \item Participação em eventos acadêmicos e tecnológicos
  \item Promoção ativa do uso de sistemas Linux e outros projetos de \textit{software} livre e código aberto
\end{conduta}

\subsection{Uso de Infraestrutura}

Os recursos e canais do grupo devem ser utilizados exclusivamente para fins acadêmicos, técnicos e colaborativos, 
em conformidade com os princípios do Manifesto e as normas do Código de Conduta.

\subsection{Propriedade Intelectual}

Projetos desenvolvidos no âmbito do grupo deverão, obrigatoriamente, adotar licenças de \textit{software} livre e código aberto, 
garantindo liberdade de uso, modificação e distribuição.

\subsection{Membros Fundadores}

Os signatários desta Carta de Fundação são reconhecidos como membros fundadores do GULECD/UFPR.

%--- Seção 6: Disposição Final ----------------------------------------------
\section{Disposição Final}

A presente Carta de Fundação entra em vigor a partir de sua publicação e estabelece as bases institucionais, 
operacionais e éticas do Grupo de Usuários Linux e Software Livre do curso de Estatística e Ciência de Dados da UFPR, consolidando seu compromisso 
com a construção de uma comunidade técnica de excelência, orientada pela liberdade, pelo rigor científico e pela colaboração aberta.

A presente Carta de Fundação constitui registro histórico do ato de criação do GULECD/UFPR, não estando sujeita a alterações após sua formalização.

Os demais documentos institucionais do grupo são considerados automaticamente aceitos pelos signatários desta carta.

Eventuais evoluções institucionais serão registradas em documentos complementares, sem alteração do presente instrumento.

%--- Espaço para assinaturas ------------------------------------------------
\vspace{1cm}
\begin{center}
	\rule{0.45\textwidth}{0.4pt}\hfill\rule{0.45\textwidth}{0.4pt}\\[0.3em]
	\small\color{gulecdGray}
	Assinatura da Coordenação Fundadora
\end{center}

\vspace{2cm}
\begin{center}
	\rule{0.45\textwidth}{0.4pt}\\[0.3em]
	\small\color{gulecdGray}
	Coordenador Fundador
\end{center}

\vspace{2cm}
\begin{center}
	\rule{0.45\textwidth}{0.4pt}\\[0.3em]
	\small\color{gulecdGray}
	Vice-Coordenador Fundador
\end{center}

\vspace{4cm}
\begin{center}
	\rule{0.45\textwidth}{0.4pt}\hfill\rule{0.45\textwidth}{0.4pt}\\[0.3em]
	\small\color{gulecdGray}
	Assinatura do Núcleo de Coordenadores Fundadores
\end{center}



%============================================================================
\end{document}
%============================================================================
